\section{Morava descent and the actual height-two test}\label{sec:homotopy}

The ordinary coefficient theorem computes four cohomology lines. We now
place that calculation inside a spectral sequence of actual spectra.
The goal of this section is to prove convergence after the height-two test
and identify its completed terms as inverse limits of ordinary cooperation
quotients. Section~\ref{sec:jets} will compare those quotients through every
power of the height-two parameter.

Let \(E_i=E_i(k)\) be Morava \(E\)-theory for the chosen \(\Gamma_i\),
for \(i=2,3\) \citep{DH04,HS99}. Recall the exact test \(\mathcal T\) of
Equation~\eqref{eq:generic-test} and the cotangent periodicity line
\(\Omega=\omega/p\) over \(\mathcal O_x=k[[x]]\). Its tensor powers include
inverse powers. Unless a relative tensor product is displayed explicitly,
\(\wedge\) denotes the ordinary smash product of spectra.

\subsection{Spherical constants and unextended descent}

\begin{lemma}[Spherical constants]\label{lem:spherical-constants}
There is a finite \'{e}tale $S$-algebra $S_k$ with
$\pi_0S_k=W(k)$.  If $d=[k:\Fp]$, then
\begin{equation}\label{eq:spherical-constants}
  \pi_*S_k=\pi_*S\otimes_{\Zp}W(k),\qquad
  S_k\simeq S^{\oplus d}
  \quad\text{as $S$-modules}.
\end{equation}
The coefficient inclusions give canonical maps $S_k\to E_i(k)$, and
stabilizer automorphisms fixing $k$ are $S_k$-linear.  Set
\begin{equation}\label{eq:test-spectra}
  Z_k=L_{K(3)}S_k,\qquad
  Y=L_{K(2)}(E_2\wedge Z_k).
\end{equation}
The algebra $E_3(k)$ is descendable over $Z_k$ in the $K(3)$-local
category.  For its Amitsur object
\begin{equation}\label{eq:morava-amitsur}
  C^s=\underbrace{E_3\otimes^{K(3)}_{Z_k}\cdots
                    \otimes^{K(3)}_{Z_k}E_3}_{s+1\text{ factors}},
\end{equation}
completed Morava cooperations identify
\begin{equation}\label{eq:unextended-cooperations}
  \pi_*C^s=\operatorname{Map}_{\mathrm{cts}}(P^s,E_{3*}),
\end{equation}
with the adic coefficient topology and the usual evaluation and
stabilizer-action faces.
\end{lemma}

\begin{proof}
\emph{Finite constants and the selected Morava factor.}
The classification of \'{e}tale algebras over a ring spectrum
\cite[Definition~2.31 and Theorem~2.32]{Mathew16} lifts
$W(k)/\Zp$ and gives the first identity in
\eqref{eq:spherical-constants}.  Represent a $\Zp$-basis of $W(k)$ by
elements of $\pi_0S_k$.  The resulting map $S^{\oplus d}\to S_k$ is an
isomorphism on every homotopy group, proving the second assertion.
The \'{e}tale mapping property stated immediately after Theorem~2.32 in
\citet{Mathew16} lifts the coefficient
inclusion into $\pi_0E_i(k)$ uniquely and supplies its compatibility with
automorphisms.  In particular,
$Z_k\simeq (L_{K(3)}S)^{\oplus d}$ as an underlying spectrum.

Here is the constant-field argument for descent.  Put
$k_0=\mathbb F_{p^3}$, $E^{(0)}=E_3(k_0)$, and $T=L_{K(3)}S$.
The algebra $E^{(0)}$ is descendable over $T$
\cite[Theorem~4.18 and Proposition~10.10]{Mathew16}.  The images of
$S$ and of the sphere in this category agree, since completion of the
sphere is a mod-$p$ equivalence.  Base change preserves descendability
\cite[Corollary~3.21]{Mathew16}; hence
\[
  \widetilde E=Z_k\otimes^{K(3)}_T E^{(0)}
\]
is descendable over $Z_k$.  Finite \'{e}tale scalar extension and
\[
  W(k_0)\otimes_{\Zp}W(k)
    =\prod_{\alpha:k_0\hookrightarrow k}W(k)
\]
give a product decomposition
\[
  \widetilde E\simeq\prod_{\alpha:k_0\hookrightarrow k}E_\alpha,
  \qquad E_\alpha\simeq E_3(k).
\]
Each factor realizes the universal deformation after the indicated
residue-field extension.  Frobenius in the extended stabilizer of
$E^{(0)}$ induces a $Z_k$-linear automorphism of $\widetilde E$ that
cycles its three factors.  Consequently $\widetilde E$ and any selected
$E_\alpha$ generate the same thick tensor ideal of $Z_k$-modules:
one direction uses the product projections, and the other uses the
equivalences of the three factors.  This proves descendability of the
selected $E_3(k)$.  In particular, no averaging over a group of order
three is involved.

\smallskip\noindent\emph{The unextended cooperation maps.}
The factor selection also determines the group and the maps in
\eqref{eq:unextended-cooperations}.  Write
$G=P\rtimes\operatorname{Gal}(k_0/\Fp)$.  The completed Morava
cooperation theorem
\cite[Proposition~2.2 and Equations~(2.3) and~(2.5)]{DH04}, followed by
the finite scalar extension above, gives
\[
  \pi_*\bigl(\widetilde E^{\otimes^{K(3)}_{Z_k}(s+1)}\bigr)
     =\operatorname{Map}_{\mathrm{cts}}(G^s,\widetilde E_*).
\]
Finite free scalar extension commutes with these continuous-function
modules.  We spell out the conversion from the cited inverse-action
coordinates. Write $\mu$ for iterated multiplication. In the coordinates
$h_1,\ldots,h_s$ of \citet[Equation~(2.5)]{DH04}, evaluation is
\[
 \mu\circ(h_1^{-1}\otimes\cdots\otimes h_s^{-1}\otimes1),
\]
where $1$ denotes the identity action.
For $a_r=g_1\cdots g_r$ and $a_0=1$, put
$h_i=a_{i-1}^{-1}a_s=g_i\cdots g_s$ and postcompose by $a_s$.
Multiplicativity of the action gives the identity of evaluation composites
\begin{equation}\label{eq:amitsur-coordinate-conversion}
 a_s\circ\mu\circ(h_1^{-1}\otimes\cdots\otimes h_s^{-1}\otimes1)
 =\mu\circ(1\otimes a_1\otimes\cdots\otimes a_s).
\end{equation}
This is a continuous invertible coordinate change, with
$g_i=h_i h_{i+1}^{-1}$ and $h_{s+1}=1$; it moves the coefficient action
from the last factor to the first. A pure tensor is therefore sent to
\[
 (g_1,\ldots,g_s)\longmapsto
 b_0\,g_1(b_1)\,(g_1g_2)(b_2)\cdots(g_1\cdots g_s)(b_s).
\]
The same identity of evaluation composites checks the bar operations.
For a cochain $f$, unit insertion in the first slot gives
$g_1 f(g_2,\ldots,g_{s+1})$;
an interior insertion gives
$f(g_1,\ldots,g_i g_{i+1},\ldots,g_{s+1})$; and insertion in the last
slot gives $f(g_1,\ldots,g_s)$. Multiplication of neighboring slots gives
insertion of the identity among the group variables. These are the actual
cofaces and codegeneracies after the coordinate change.
Selecting the same factor idempotent in all $s+1$ slots first selects
that value factor and then requires every cumulative product
$g_1\cdots g_r$ to preserve it.  The stabilizer of the factor is $P$;
these conditions are equivalent to $g_r\in P$ for every $r$.
This yields \eqref{eq:unextended-cooperations}.  The selection commutes
with multiplication, units, and evaluation, so the resulting faces are
the actual bar faces, not merely maps with the same groups of
coefficients.
\end{proof}

\subsection{Convergence after the height-two test}

The test \(\mathcal T(X)=L_{K(2)}(E_2\wedge X)/p\) takes values in
\(E_2\)-modules equipped with the semilinear action induced by \(J\) on
\(E_2\). Its exactness is the reason a finite nilpotence bound on the
Amitsur error tower survives this change of height.

\begin{lemma}[Partial totalizations as finite cubes]\label{lem:finite-totalizations}
Let $X^\bullet$ be a cosimplicial object in a stable $\infty$-category.
For $r\geq0$, let $\mathcal P_0([r])$ be the poset of nonempty subsets of
$[r]=\{0,\ldots,r\}$.  The functor
\[
 g_r:\mathcal P_0([r])\longrightarrow\Delta_{\leq r},\qquad
 S\longmapsto[|S|-1],
\]
sends an inclusion to the injection recording the positions of its ordered
elements.  Here $\Delta_{\leq r}$ is the full subcategory on
$[0],\ldots,[r]$, including all codegeneracies.  Restriction induces
\begin{equation}\label{eq:finite-totalization-cube}
 \Tot_r X^\bullet=\lim_{\Delta_{\leq r}}X^\bullet
 \simeq\lim_{S\in\mathcal P_0([r])}X^{|S|-1}.
\end{equation}
These equivalences respect the totalization tower and the augmentation
when $X^\bullet$ is augmented.  Every exact functor between stable
$\infty$-categories therefore preserves these partial totalizations as
an augmented tower.
\end{lemma}

\begin{proof}
The assertion about the indexing functor is the left-cofinality theorem
of \citet[Definitions~6.2--6.3 and Theorem~6.7]{Sinha09}.
We include a proof that records the role of codegeneracies.  For
$0\leq m\leq r$, the comma category $(g_r\downarrow[m])$ consists of
nonempty subsets $S\subseteq[r]$ with weakly increasing labelings
$S\to[m]$; its arrows extend labelings.  It is the nonempty face poset of
the finite simplicial complex $K_{r,m}$ with simplices
\[
 \{(i_0,j_0),\ldots,(i_q,j_q)\},\qquad
 i_0<\cdots<i_q,\quad j_0\leq\cdots\leq j_q.
\]
This complex is flag.  In a flag complex one may delete a vertex $v$
if a neighboring vertex $w$ is adjacent to every other neighbor of $v$:
sending $v$ to $w$ and fixing the other vertices gives a simplicial
retraction whose composite with the inclusion is contiguous to the identity.

For $j=m,m-1,\ldots,1$, delete $(i,j)$ for $i=0,\ldots,j-1$ in that
order, using $w=(i+1,j)$.  This vertex exists since $j\leq m\leq r$.
The only neighbors of $(i,j)$ possibly not adjacent to $w$ are
$(i+1,k)$ with $k>j$; they have already been deleted at the earlier
value $k$.  Next, for $j=0,1,\ldots,m$, delete $(i,j)$ for
$i=r,r-1,\ldots,j+1$, using $w=(i-1,j)$.  The only possible exceptions
are now $(i-1,k)$ with $k<j$; they were deleted at the earlier value
$k$.  Each chosen $w$ is still present.  The vertices left are
$(0,0),\ldots,(m,m)$, which span a simplex.  Thus $K_{r,m}$, and hence
the nerve of $(g_r\downarrow[m])$, is contractible.  The weak label
inequalities here include noninjective simplex maps; this is the check
for the full truncation.

The opposite form of the $\infty$-categorical cofinality criterion
\cite[Theorem~4.1.3.1]{Lurie09} makes $g_r$ initial for limits.
The dual of \cite[Proposition~4.1.1.8]{Lurie09} gives
\eqref{eq:finite-totalization-cube}.  The poset on its right has a
finite nerve, since strict chains of its subsets have length at most $r$.
For the inclusions of initial segments, the squares formed by $g_r$ and
$g_{r+1}$ commute strictly.  The displayed equivalences, being restriction
maps on cones, consequently commute coherently with the tower maps.
Adding the empty subset, sent to $[-1]$ in the augmented simplex category,
shows the same compatibility with the augmentation.  Finally, an exact
functor preserves the finite limit on the right, naturally in these diagrams.
Applying the comparison in source and target proves the last assertion.
\end{proof}


\begin{proposition}[The tested descent spectral sequence]
\label{prop:descent-test}
The natural augmentation induces an equivalence
\begin{equation}\label{eq:tested-totalization}
  Y/p\simeq\Tot\mathcal T(C^\bullet).
\end{equation}
The corresponding spectral sequence
\begin{equation}\label{eq:tested-spectral-sequence}
  \mathsf E_2^{s,t}
     =H^s\bigl(\pi_t\mathcal T(C^\bullet)\bigr)
       \Longrightarrow\pi_{t-s}(Y/p)
\end{equation}
is strongly convergent, with a horizontal vanishing line at a finite
page and a finite filtration on its abutment.  Its differentials are
$k[[x]]$-linear and $J$-semilinear-equivariant.

There is a compatible diagonal summand, denoted by a subscript
$\Delta$, on every page and on the abutment filtration: it is the
summand on which the two copies of $k$, from $E_2$ and $S_k$, agree.
The same assertions hold on that summand.
\end{proposition}

\begin{proof}
\emph{The Amitsur error tower.}
Let $T_r=\Tot_r C^\bullet$.  In the stable category of $Z_k$-modules
in $K(3)$-local spectra, form the external tensor cube
\[
 Q_r(S)=\bigotimes_{i=0}^{r} Z_i(S),\qquad
 Z_i(S)=
 \begin{cases}E_3,&i\in S,\\ Z_k,&i\notin S,\end{cases}
 \qquad S\subseteq[r],
\]
where the tensors are the $K(3)$-local relative tensors over $Z_k$.
An edge applies the unit $Z_k\to E_3$ in its ordered slot.  More generally,
the Amitsur map for a weakly increasing simplex map multiplies the factors
in each ordered fiber and inserts a unit for an empty fiber; its
codegeneracies are the multiplication maps.  The restriction along $g_r$
uses precisely the unit insertions.  The associativity and unit equivalences
therefore identify it with the punctured cube $Q_r|_{\mathcal P_0([r])}$ as
a coherent diagram, and the empty vertex is the actual augmentation $Z_k$.
Lemma~\ref{lem:finite-totalizations} gives
\[
 T_r\simeq\lim_{\varnothing\ne S\subseteq[r]}Q_r(S)
\]
as an augmented tower.  For $r=0$ the limit is $E_3$; for $r=1$ it is the
homotopy pullback of
$E_3\to E_3\otimes^{K(3)}_{Z_k}E_3\leftarrow E_3$, with ordered maps
$b\mapsto b\otimes1$ and $b\mapsto1\otimes b$.

Put $I=\operatorname{fib}(Z_k\to E_3)$.  Iterating fibers in the tensor
cube, using exactness of the tensor in each factor, gives
\begin{equation}\label{eq:amitsur-error-fiber}
 \operatorname{fib}(Z_k\to T_r)
       \simeq I^{\otimes(r+1)}.
\end{equation}
Indeed, after separating the last coordinate, the total fibers of the two
faces are $I^{\otimes r}$ and $E_3\otimes I^{\otimes r}$; their fiber is
$I^{\otimes(r+1)}$.  The case $r=0$ is the definition of $I$.  Naturality
of this iterated-fiber construction identifies a tower transition with
applying $I\to Z_k$ in the last slot.

Take the cofiber of the augmentation in the category of towers,
\[
  Z_k\longrightarrow T_r\longrightarrow U_r.
\]
Thus $U_r\simeq\Sigma I^{\otimes(r+1)}$ compatibly with transitions.
Multiplication retracts
$E_3\simeq E_3\otimes Z_k\to E_3\otimes E_3$.  Applying that retraction
to the defining fiber nullhomotopy shows that
$E_3\otimes I\to E_3$ is null.  Therefore each transition
$U_r\to U_{r-1}$ becomes null after tensoring with $E_3$.
Descendability supplies a bound on the length of any composite of such maps
\cite[Proposition~3.27 and Corollary~4.4]{Mathew16}.  Thus there is
$N$ such that $U_{r+N}\to U_r$ is null for every $r$.

\smallskip\noindent\emph{Passing the tower through the exact test.}
This uniform bound is the descent input that will survive the change
of height. The test need only preserve the finite limits defining the
partial totalizations.
Lemma~\ref{lem:finite-totalizations} identifies
$\mathcal T(T_r)$ with $\Tot_r\mathcal T(C^\bullet)$ as augmented towers:
exactness moves $\mathcal T$ through the finite punctured-cube limit.
Here $\mathcal T$ is applied to the already formed source cube; no
monoidal property of the test is needed.  Exactness then gives the cofiber tower
\[
  \mathcal T(Z_k)\longrightarrow
       \Tot_r\mathcal T(C^\bullet)\longrightarrow\mathcal T(U_r),
\]
and every composite $\mathcal T(U_{r+N})\to\mathcal T(U_r)$ is still
null.  Increase $N$ to at least one if necessary, and write
\[
 X=\mathcal T(Z_k),\qquad
 V_s=\Tot_s\mathcal T(C^\bullet),\qquad W_s=\mathcal T(U_s).
\]
The inverse limit of $W_s$ is zero: applying mapping spectra gives
pro-zero homotopy groups and zero $\lim^1$.  Taking the inverse limit
of the displayed cofiber sequences gives $X\simeq\lim_sV_s$ and proves
\eqref{eq:tested-totalization}.  The tower $\{V_s\}$ is strongly
constant in the sense of \cite[Definition~4.2]{Mathew16}, and
\cite[Proposition~4.3]{Mathew16} supplies a horizontal vanishing line
at a finite page.

\smallskip\noindent\emph{The filtration of actual homotopy.}
We derive the abutment filtration and strong convergence directly from
the same nilpotence bound. This will let us use a surviving term as a
subquotient of the tested homotopy group.
For any homotopy degree $q$ and $u\ge s+N$, the cofiber triangles give
\[
 \operatorname{im}(\pi_qV_u\longrightarrow\pi_qV_s)
 =\operatorname{im}(\pi_qX\longrightarrow\pi_qV_s)=:M_s^q.
\]
Indeed, the image of a class from $\pi_qV_u$ in $\pi_qW_s$ is zero,
because it factors through the null transition $W_u\to W_s$; exactness
then puts its image in the image of $\pi_qX$.  Conversely, the augmentation
to $V_s$ factors through every $V_u$.  Hence each homotopy tower is
Mittag--Leffler, and its $\lim^1$ vanishes.  The Milnor exact sequence
therefore identifies $\pi_qX$ with $\lim_s\pi_qV_s$.

Extend $V_s=0$ for $s<0$ and give this actual homotopy group the filtration
\[
 F^s\pi_qX=\ker(\pi_qX\longrightarrow\pi_qV_{s-1})\qquad(s\ge0).
\]
Here $F^0\pi_qX=\pi_qX$.  Exactness identifies $F^{s+1}\pi_qX$ with
the image of the boundary $\delta_s:\pi_{q+1}W_s\to\pi_qX$.
Naturality gives
$\delta_N=\delta_0\circ\pi_{q+1}(W_N\to W_0)=0$, so
$F^{N+1}\pi_qX=0$ for every $q$.

Use total-degree grading
$\widetilde E_r^{s,q}=\mathsf E_r^{s,q+s}$ in the exact couple of the tower.
Its first couple has
\[
 D_1^{s,q}=\pi_qV_s,\qquad
 \widetilde E_1^{s,q}=
   \pi_q\operatorname{fib}(V_s\longrightarrow V_{s-1}).
\]
The long exact fiber sequences give maps $i_1:D_1^{s,q}\to D_1^{s-1,q}$,
$j_1:D_1^{s,q}\to\widetilde E_1^{s+1,q-1}$, and
$\epsilon_1:\widetilde E_1^{s,q}\to D_1^{s,q}$.
Deriving the couple replaces $D$ by the image of $i$.  Induction therefore
gives
\[
 D_r^{s,q}=\operatorname{im}(\pi_qV_{s+r-1}\longrightarrow\pi_qV_s),
\]
with maps
\[
 i_r:D_r^{s,q}\to D_r^{s-1,q},\qquad
 j_r:D_r^{s,q}\to\widetilde E_r^{s+r,q-1},\qquad
 \epsilon_r:\widetilde E_r^{s,q}\to D_r^{s,q},
\]
and differential $\widetilde d_r=j_r\epsilon_r$.  In the original grading
this has bidegree $(r,r-1)$.

Put $M_s^q=0$ for $s<0$.  For $r\ge N+1$, the stable-image equality
says $D_r^{s,q}=M_s^q$ for every $s$.  The transitions $M_s^q\to M_{s-1}^q$
are surjective because both images come from $\pi_qX$.  Exactness of the
derived couple now makes every $j_r$ zero and identifies
\[
 \widetilde E_r^{s,q}\xrightarrow[\epsilon_r]{\ \sim\ }
       \ker(M_s^q\longrightarrow M_{s-1}^q).
\]
Indeed, in the exact segment
$D_r^{s-r,q+1}\to\widetilde E_r^{s,q}\to D_r^{s,q}\to D_r^{s-1,q}$,
the first arrow is zero since the preceding $i_r$ is surjective.
Thus the pages stabilize from $r=N+1$.  The augmentation maps
$F^s\pi_qX$ onto the displayed kernel with kernel $F^{s+1}\pi_qX$,
so in the original grading
\[
 \mathsf E_\infty^{s,t}
    \cong F^s\pi_{t-s}X/F^{s+1}\pi_{t-s}X\qquad(s\ge0).
\]
This is an exhaustive, separated, and complete finite filtration of the
actual group $\pi_{t-s}(Y/p)$: it starts at $F^0$ and ends at $F^{N+1}=0$.
The stable page consequently vanishes for $s\ge N+1$, and the Milnor term
above is zero.  This proves strong convergence with the asserted actual
abutment filtration.  The argument uses only exactness of $\mathcal T$
through the finite cube, not preservation of arbitrary inverse limits.

\smallskip\noindent\emph{Coefficient actions and the diagonal summand.}
The maps in the tested cosimplicial object are $E_2$-linear and
compatible with the action of $J$.  Lemma~\ref{lem:mod-p-nullity}
makes their homotopy groups, and hence the spectral sequence,
$k[[x]]$-linear.  The second coefficient action comes from $S_k$ on
$C^\bullet$ and commutes with the first.  After reduction modulo $p$
the two actions give an action of
\[
  k\otimes_{\Fp}k\simeq
          \prod_{\sigma\in\operatorname{Gal}(k/\Fp)}k.
\]
Its idempotent for $\sigma=\id$ is fixed by $P$ and $J$.  It commutes
with the tower maps, all differentials, and the maps defining the
abutment filtration.  Taking its image proves the diagonal assertions.
This construction only needs an idempotent on mod-$p$ homotopy groups;
it does not identify the two spherical scalar actions on an ordinary
smash product before reduction.
\end{proof}

\subsection{Ordinary cooperations and their residue topology}

For $m\geq1$, put $A_m=k[x]/(x^m)$ and
$\omega_m=\Omega\otimes_{k[[x]]}A_m$.  All the quotients below are
formed on the height-two coefficient side.  Let
\begin{equation}\label{eq:jet-cooperation-definition}
  D_m^s=
  \left(\pi_0(E_2\wedge C^s)/(p,x^m)\right)_\Delta.
\end{equation}
The cofaces and degeneracies make $D_m^\bullet$ a cosimplicial
$A_m$-algebra; we use the same notation for its associated cochain
complex.  The action of $J$ on the algebraic quotient is well defined
because it preserves $(p,x^m)$.

\begin{lemma}[Flat coefficients and the residue algebra]
\label{lem:ordinary-cooperations}
The ordinary, unlocalized smash product has even homotopy and
\begin{equation}\label{eq:ordinary-kunneth}
  \pi_*(E_2\wedge C^s)
    =(E_2)_*(E_3)\otimes_{E_{3*}}
          \operatorname{Map}_{\mathrm{cts}}(P^s,E_{3*}).
\end{equation}
This module is flat over $E_{2*}$.  In particular, $D_m^s$ is flat
over $A_m$, and reduction by $(p,x^m)$ creates no odd homotopy.
There are natural identifications
\begin{equation}\label{eq:residue-cochains}
  D_1^s=
  \operatorname{Map}_{\mathrm{cts}}(P^s,K)\otimes_K L
    =\colim_U\operatorname{Map}_{\mathrm{cts}}(P^s,L^U),
\end{equation}
compatible with the bar maps and the marking actions.
\end{lemma}

\begin{proof}
\emph{Evenness and flatness before localization.}
The Landweber exact cooperation formula expresses $(E_2)_*(E_3)$
as the two-sided base change of the formal-group isomorphism Hopf
algebroid.  It is even and flat over either coefficient ring.  These
are the usual ordinary Landweber cooperations; see
\cite[Theorem~2.7 and Propositions~2.12 and~2.16]{HS99}.
For flatness over both coefficient rings, apply Proposition~2.16
with the two Morava spectra in each order and use smash symmetry.
We justify this ordinary tensor calculation by a finite resolution.
By \cite[Proposition~2.12]{HS99}, $E_2$ is evenly generated;
\cite[Proposition~2.16]{HS99}, applied with ring spectrum $E_3$,
then gives a graded projective resolution
\[
  0\longrightarrow P_1\longrightarrow P_0
       \longrightarrow\pi_*(E_2\wedge E_3)\longrightarrow0.
\]
The projectives may be taken in even degrees. Realize $P_0$ as a
retract $Q_0$ of a wedge of even suspensions of $E_3$, and realize
the surjection by an $E_3$-module map $Q_0\to E_2\wedge E_3$.
Its fiber has homotopy $P_1$, so realizing that projective similarly
identifies the fiber with a projective $E_3$-module $Q_1$. Thus
$Q_1\to Q_0\to E_2\wedge E_3$ is a finite cofiber sequence.
Tensor this sequence over $E_3$ with $C^s$. For each $Q_i$, its
homotopy is the algebraic tensor product
$P_i\otimes_{E_{3*}}\pi_*C^s$, since $Q_i$ is a retract of a wedge
of suspensions of $E_3$. Flatness of $\pi_*(E_2\wedge E_3)$ over
$E_{3*}$ makes the map between these two tensor products injective.
We also have
\[
  E_2\wedge C^s\simeq(E_2\wedge E_3)\otimes_{E_3}C^s.
\]
The long exact homotopy sequence therefore gives
\eqref{eq:ordinary-kunneth} and evenness. This finite
argument supplies the tensor calculation without a convergence
assumption on an unbounded periodic K\"unneth spectral sequence.

We also need flatness of the second factor in
\eqref{eq:ordinary-kunneth}.  The compact analytic space $P^s$ has a
countable cofinal system of finite clopen partitions.  The modules of
functions constant on these partitions are finite free over $E_{3,0}$.
A refinement map is a split inclusion with free cokernel, so their
union is free.  Its maximal-ideal completion is
$\operatorname{Map}_{\mathrm{cts}}(P^s,E_{3,0})$: reduction modulo any
power of the maximal ideal is a locally constant function, and
compatible reductions recover a continuous function.  Completion of
a flat module over a Noetherian ring is flat.  Thus this
continuous-function module is flat over $E_{3,0}$.

To see flatness of the full tensor product over $E_{2*}$, first
tensor an exact sequence of $E_{2*}$-modules with
$(E_2)_*(E_3)$, and then tensor the resulting exact sequence of
$E_{3*}$-modules with the continuous-function module.  Both operations
are exact.  It follows that $p,x^m$ is a regular quotient calculation
on \eqref{eq:ordinary-kunneth}.  The same holds after taking the
diagonal idempotent.  Flatness over $A_m$ follows by base change.

\smallskip\noindent\emph{The residue algebra as a marking algebra.}
Choose complex orientations. Quillen's universality theorem identifies
the complex-cobordism formal group with the universal formal group
\cite[Section~3, Theorem~2]{Quillen69}. We use its strict-isomorphism
Hopf algebroid as in \cite[Section~3.1]{Powell12}, with the inverse
series when the direction of the isomorphism is reversed.
The ordinary cooperation identity is
\[
 (E_2)_*E_3=(E_2)_*\otimes_{MU_*}MU_*MU
                         \otimes_{MU_*}(E_3)_*,
 \qquad MU_*MU=MU_*[b_1,b_2,\ldots],
\]
with the two unit maps of the formal-group Hopf algebroid.
Both tensor products are algebraic.  This follows by applying
Landweber exactness \cite[Theorem~2.7]{HS99} first to $(E_2)_*(E_3)$ and then, using symmetry,
to $MU_*(E_3)=(E_3)_*(MU)$.

Choose periodicity generators $u_i\in\pi_2E_i$.  If
$b(Z)=Z+\sum_{r\geq1}b_rZ^{r+1}$ is the strict graded
isomorphism from the right formal law to the left, its degree-zero
form is
\[
          f(X)=u_2b(u_3^{-1}X),\qquad f'(0)=c=u_2/u_3.
\]
Conversely, any degree-zero isomorphism with invertible derivative
$c$ recovers $u_3=u_2/c$ and the strict isomorphism
$u_2^{-1}f(u_3Z)$.  The periodicity ratio therefore supplies exactly
the arbitrary linear coefficient of a formal-law isomorphism.

Reduce the left coefficients by $(p,x)$ and select the diagonal
copy of $k$.  The left law is $\Gamma_2$, with
$[p]_{\Gamma_2}(X)=X^{p^2}$.  For the right law $F_3$ over
$A=k[[a,t]]$, compare coefficients in
\[
          f([p]_{F_3}(X))=f(X)^{p^2}.
\]
The coefficient of $X^p$ gives $ca=0$, hence $a=0$.
The coefficient of $X^{p^2}$ then gives
$ct=c^{p^2}$, so $t=c^{p^2-1}$ is invertible.
The right coefficient ring is consequently
$A/(a)[1/t]=K$.  The remaining coefficient relations classify
exactly the isomorphisms from the connected formal law of
$\mathcal G_K$ to $\Gamma_2$.  By the construction of the
connected-marking torsor, their algebra is
$L=\colim_U L^U$: each element uses finitely many marking
coefficients and lies in one finite \'{e}tale marking algebra.
This is the algebraic union, with no additional completion.

\smallskip\noindent\emph{Continuous functions at a finite marking stage.}
On the continuous-function factor, reduction modulo $(p,a)$ gives
the continuous functions into $k[[t]]$.  Indeed coefficientwise
lifts give surjectivity, and division of an element in the kernel by
$p$ or $a$ is continuous with a bounded shift of adic order.
After inverting $t$ the result is
$\operatorname{Map}_{\mathrm{cts}}(P^s,K)$: the compact image of a
continuous $K$-valued function has bounded pole order, giving one
common power of $t$ as denominator.

Finally, every finite \'{e}tale marking stage is a finite-dimensional
$K$-algebra with its usual finite-dimensional topology.  A choice of
$K$-basis therefore identifies continuous functions into that stage
with its tensor product with
$\operatorname{Map}_{\mathrm{cts}}(P^s,K)$.  Passing to the algebraic
union proves \eqref{eq:residue-cochains}.  Every cooperation element
uses only finitely many marking coefficients, so this is a union of
complete finite stages, with no completion of the infinite extension.
Evaluation and transport of the height-three law are exactly the
cooperation faces, giving the asserted compatibility.
\end{proof}

\begin{lemma}[The completed terms]\label{lem:completed-test-terms}
The diagonal homotopy cochains of $\mathcal T(C^\bullet)$ are zero
in odd internal degrees.  In degree $2j$ they are naturally
\begin{equation}\label{eq:completed-homotopy-cochains}
  \left(\lim_m D_m^\bullet\right)
             \otimes_{k[[x]]}\Omega^{\otimes j}.
\end{equation}
The maps to the finite quotients on homotopy groups are the ordinary
reduction maps.
\end{lemma}

\begin{proof}
For $E_2$-modules, $K(2)$-localization is derived completion at
$(p,x)$: it is the inverse limit of the successive regular-ideal
quotients \cite[Proposition~2.2 and Theorem~2.3]{Hovey04}; compare
\cite[Proposition~7.10(e)]{HS99}.
Here is the comparison for the finite extension of constants. In
Section~\ref{sec:setup} the height-two parameter was chosen already over
$\mathbb F_{p^2}=\mathbb F_9$, and $E_2(k)$ uses its scalar-extended
universal deformation. The map
$E_2(\mathbb F_9)\to E_2(k)$ therefore sends the chosen regular system
$(p,x)$ to the same named system. Apply the standard-coefficient theorem
to an $E_2(k)$-module restricted along this map, using the regular ideals
$(p^r,x^r)$, whose intersection is zero. The underlying spectrum and
its $K(2)$-localization do not change under restriction. Multiplication
by $p^r$ and $x^r$ is the same map on that spectrum, so its derived
quotients and their transition maps are also the same. This proves the
same completion description for $E_2(k)$-modules.
Exactness first permits reduction
modulo $p$ before localization.  On this reduction $p$ acts null by
Lemma~\ref{lem:mod-p-nullity}; hence derived $(p,x)$-completion is
derived $x$-completion.  Equivalently, tensoring the usual completion
tower with the finite $E_2$-module $E_2/p$ commutes with its inverse
limit, and the extra $p$-power tower has pro-zero error.  Thus, for
$M=E_2\wedge C^s$,
\[
  L_{K(2)}M/p\simeq\lim_m M/(p,x^m).
\]
Lemma~\ref{lem:ordinary-cooperations} computes every term on the
right by the ordinary even quotient of \eqref{eq:ordinary-kunneth}.
The homotopy transition maps are surjective, so the Milnor
$\lim^1$ term vanishes.  Periodicity on the height-two factor then
gives \eqref{eq:completed-homotopy-cochains}.  Its line tensor
commutes with the limit because it is finite free as an underlying
$k[[x]]$-module.  These constructions are natural in all bar maps.
\end{proof}


We have now identified the homotopy of each actual completed term and its
reduction maps. The remaining comparison must respect those maps at every
finite power of \(x\); a calculation only at \(x=0\) would leave the
generic semilinear action undetermined.
